\section{The Method of Containers}

\subsection{Some nice problems}
There are many reasons to understand the method of hypergraph
containers. Here are some cool problems we may be interested
in.
\begin{problem}
	\label{containers:prob:extremalK3}
	(Extremal questions in $G(n,p)$). We define
	\[
		\text{ex}(G,H) := \max \{ e(G'): G' \subset G,\; G \text{ is }
		H-\text{free}\}.
	\]
	Note that $\text{ex}(n,H) = \text{ex}(K_n,H)$. What can we say about
	$\text{ex}(G(n,p),K_3)$?
\end{problem}

Some easy bounds can be found. Given any sampling of $G(n,p)$, we can
always find a bipartite graph with half the edges. Similarly
another strategy is just to kill off triangles by removing an edge
from each. This already gives
\[
	\text{ex}(G(n,p), K_3) \geq
	\begin{cases}
		e(G(n,p))/2\\
		e(G(n,p)) - k_3(G(n,p))
	\end{cases}
\]
by finding the best regime for $p$ in each we get
\[
	\text{ex}(G(n,p), K_3) \geq
	\begin{cases}
		pn^2/4 \quad \text{for any}\,p \in [0,1]\\
		(1 - o(1))pn^2/2 \quad \text{for}\, p \ll \frac{1}{\sqrt{n}}
	\end{cases}
\]
it turns out this is sharp.

\begin{problem}
	When does $G(n,p) \to K_3$? Define
	\[
		p^*(n) = \inf_{p \in [0,1]} \{p : \Pr(G(n,p) \to K_3) \geq 1/2\}.
	\]
	can we find this threshold?
\end{problem}
We know that if $G(n,p)$ contains a $K_6$ we are done by Ramsey so
$p^* \lesssim n^{-2/5}$. But is this it?
Should avoiding $K_6$ be enough or is there some deeper more complex
structure - the answer is obviously yes.

\begin{problem}
	Let $A \subset [n]_p$ with $|A| \geq \delta np$, does $A$ contain a $3$-AP?
\end{problem}
Doing deletion we get that it is true if  $AP_3([n]_p) = o(np)$. We
know $AP_3([n]_p) \eqsim n^2p^3$, so if $p \ll 1/\sqrt{n}$
we get the result. What happens if $p \gg 1/\sqrt{n}$?

Very related to \ref{containers:prob:countingk3} is question:
\begin{problem}
	\label{containers:prob:countingk3}
	How many $K_3$-free graphs on $n$ vertices with $m$ edges are there?
\end{problem}
If we call this quantity $X(n,m)$, we know (because we can always
pick from a bipartite graph) that
\[
	X(n,m) \geq \binom{n^2/4}{m}.
\]
If $m$ is sufficiently low, we can use $G(n,p)$ for our advantage! We know that
\[
	\Pr(K_3 \not \subset G(n,p)) \eqsim \exp(-n^3p^3)
\]
if $p \ll n^{-1/2}$ by Janson and Harris inequality. Noting that $G(n,
m) \eqsim G(n,p)$ with $p = m/\binom{n}{2}$,
we get something that says
\[
	X(n,m) \eqsim \exp(-m^3/n^3)\binom{\binom{n}{2}}{m} \quad \text{if
	}\, m \ll n^{3/2}.
\]
Is the first bound optimal?

\subsubsection{Attempts at handling them}

Let's follow a naive approach for tackling
\ref{containers:prob:extremalK3}. So let's see first moment.
Let
\[Y = \#\big\{G: K_3 \not\subset G \subset G(n,p)\; \text{ with }
e(G) > m = (1/2 + \eps)p\binom{n}{2} \big\}\]
the number of such bad subgraphs of $G(n,p)$. We would like to show
$Y = 0$ with high probability if $p \gg 1/\sqrt{n}$. If we calculate
$\Ex Y$ we get
\[
	\Ex Y = \sum_{\substack{K_3 \not \subset G\\e(G) = m}} \Pr(G \subset
	G(n,p)) \geq \binom{n^2/4}{m} p^m \eqsim
	\bigg(\frac{en^2p}{4m}\bigg)^m \eqsim e^m
\]
where for the lowerbound we only considered the subsets of a single
bipartite graph. But $\Ex Y$ is \textbf{HUGE}, so can we better this?

A very interesting point made by Rob is to look where we are losing
here. Every once and while, when you look at an equation and it
explodes, it is a good idea to see why.
Maybe we could handle the bad case differently.

Here we are losing a
ton because there are many subsets of the bipartite graph, but this
seems a bit unrealistic,
a bipartite graph contains at most $1/2$ of the edges of $G(n,p)$ and
here we are asking for an $\eps$ deviation.
To be more precise, let's check another thing.

What can we say about bipartite $H$ in $G(n,p)$ with $m$ edges? Can
we find many of those?
Note that every bipartite graph $H$ is a subset of a complete
bipartite $G$ which has at most $n^2/4$ edges. So we could say
\[
	\Pr(\exists\; \text{bip.}\,H \subset G(n,p)\; \text{with } e(H) = m) =
	\sum_{\substack{K_{(n/2,n/2)} \sim G\\G \subset G(n,p)}}
	\Pr(\exists\; \text{bip.}\,H \subset G \cap G(n,p)\; \text{with } e(H) = m)
\]
It might look that we made our lives worse, but in reality this is a
way simpler thing. For such $H$ to exists, $e(G \cap G(n,p)) \geq m$, so
here we can already apply a deviation bound. Furthermore, the number
of complete bipartites is $2^{n-1}$ so applying the union bound seems
like a great idea.
We get (using \ref{lemma:expec_binomial_bounds})
\begin{align*}
	\sum_{\substack{K_{(n/2,n/2)} \sim G\\G \subset G(n,p)}}
	\Pr(\exists\; \text{bip.}\,H \subset G \cap G(n,p)\; \text{with }
	e(H) = m) &\leq
	2^n\Pr(\text{Bin}(n^2/4,p) \geq m)\\
	&\leq 2^n\exp(-\eps^2n^2p/4) \to 0
\end{align*}
if $p \gg 1/n$ so we win by miles in this bipartite case.

What we actually did here is a container theorem for bipartite graphs.
Let $\mathcal{G}$ be the family of complete bipartite graphs on $n$ vtxs,
we know that
\begin{enumerate}[label=(\alph*)]
	\item $|\mathcal{G}|$ is small $\leq 2^n$.
	\item Every $G \in \mathcal{G}$ is bipartite.
	\item For every bipartite $H \subset K_n$, there exists $G \in
		\mathcal{G}$ such that $H \subset G$. ($G$ is the container for $H$).
\end{enumerate}
This family allowed us to do way more efficient union bound. And this
is what will try to do to tackle
\ref{containers:prob:extremalK3}.

For our problem, this we need another container theorem which we will
prove later:
\begin{theorem}
	\label{trm:container-k3-free-bad}
	(Container theorem for $K_3$-free graphs). There exists a family
	$\mathcal{G}$ of graphs on $n$-vtxs, s.t
	\begin{enumerate}[label=(\alph*)]
		\item $\mathcal{G}$ is small, $|\mathcal{G}| \leq 2^{n^{3/2}(\log n)^c}$.
		\item Every $G \in \mathcal{G}$ has "few" triangles, $o(n^3)$.
		\item For every $K_3$-free graph $H$, there exists a container $G
			\in \mathcal{G}$ s.t. $H \subset G$.
	\end{enumerate}
\end{theorem}

\begin{proof}[An answer for \ref{containers:prob:extremalK3}]
	Note that we can upperbound $\Pr(\exists K_3 \not \subset H \subset
	G(n,p), e(H) = m)$ using our containers.
	How many edges can our containers have? Because they have $o(n^3)$
	triangles, by supersaturation they must have at most $(1/2 +
	o(1))\binom{n}{2}$ edges.
	So following the same calculation as before (doing the union bound
	and applying our deviation result), we get
	\[
		\Pr(\exists K_3 \not \subset H \subset G(n,p), e(H) = m) \leq
		\sum_{G \in \mathcal{G}} \Pr(e(G \cap G(n,p)) \geq m) \leq
		2^{n^{3/2}\log(n)^c}e^{-\eps pn^2} \to 0
	\]
	if $p \gg \log(n)^c/\sqrt{n}$.
\end{proof}

Next time we will try to prove Theorem \ref{trm:container-k3-free-bad},
and we will do so in greater generality. We will do it for
$3$-uniform hypergraphs and apply
it for a "triangle encoding" hypergraph.
\begin{definition}
	The triangle enconding hypergraph $\mathcal{H}_{K_3}$ is a 3-uniform
	hypergraph with $V(\mathcal{H}_{k_3}) = E(K_n)$ and
	$E(\mathcal{H}_{k_3}) = \{\{e,f,g\} : e,f,g \in E(K_n) \text{ form a
	triangle}\}$.
\end{definition}

\subsection{A general hypergraph container statement with fingerprints}
%STATE GENERAL HYPERGRAPH CONTAINERS
%PROVE THE CONTAINER THEOREM FOR TRIANGLE ENCODING GRAPHS

\begin{lemma}
	\label{lemma:hypergraph_containers}
	(Hypergraph Container Lemma) For all $c \geq 1$ and $k \in \N$,
	there exists $\delta \sim 1/ck^4$ s.t the following holds.
	Let $\mathcal{H}$ be a $k$-uniform hypergraph on $n$-vtxs and $\tau
	> 0$ satisfying
	\[
		\Delta_l(\mathcal{H}) \leq
		c\tau^{l-1}\frac{e(\mathcal{H})}{v(\mathcal{H})} = c\tau^{l-1}d
	\]
	then there exists families $\mathcal{S} \subset
	\mathcal{I}(\mathcal{H})$ and $\mathcal{C}\subset P(V(\mathcal{H}))$
	and a function $f: \mathcal{S} \to \mathcal{C}$ such that
	\begin{enumerate}[label=(\alph*)]
		\item For every $S \in \mathcal{S}$, $|S| \leq \tau n$.
		\item For every $C \in \mathcal{C}$, $|C| \leq (1-\delta)n$.
		\item For every $I \in \mathcal{I}(\mathcal{H})$, there exists
			$S(I) \in \mathcal{S}$ s.t
			\[
				S(I) \subset I \subset C(I) = f(S(I)).
			\]
	\end{enumerate}
\end{lemma}

With its full generality, we can prove the container theorem for
$K_3$ encoding graphs.
\begin{theorem}
	\label{trm:container_for_triangles}
	(Containers for triangle free graphs) For every $\eps > 0$, there
	exists a family of fingerprints $\mathcal{S}(n,\eps)$ and
	$\mathcal{G}(n,\eps)$ of subgraphs of $K_n$ satisfying the following
	properties:
	\begin{enumerate}[label=(\alph*)]
		\item $|\mathcal{G}| \leq \exp(C(\eps)n^{3/2}\log n).$
		\item For every $G \in \mathcal{G}$, $k_3(G) \leq \eps n^3.$
		\item For every triangle free graph $H$ on $n$ vtx, there exists $G
			\in \mathcal{G}$ with $H \subset G$.
	\end{enumerate}
\end{theorem}

\begin{proof}
	We will iteratively apply our lemma
	[\ref{lemma:hypergraph_containers}] until we are left with our desired family.
	What degree properties can we say about $\mathcal{H} =
	\mathcal{H}_{K_3}(K_n)$? It's simple to note that $e(\mathcal{H}) =
	\binom{n}{3}$,
	$d = \Delta_1(\mathcal{H}) = n-2$ and $\Delta_2(\mathcal{H}) \leq 1$.
 So, the first time, we can apply the lemma with $\tau = 1/\sqrt{n}$.

	Let $\mathcal{C} = \{E(K_n)\}$, while there is an element $E(G)$ in
	$\mathcal{C}$ with more than $\eps n^3$ triangles, we apply the
	container lemma on $G$, and substitute
	$G$ by its containers in $\mathcal{C}$. We know the maximum degree
	condition is the same, and because $k_3(G) \geq \eps n^3$, we know
	that the avg. degree in our
	container lemma is at least $\eps n$. So we can always repeat this
	process with $c = 1$ and $\tau = 1/\eps \sqrt{n}$.

	How many steps can we do? Well after $t$ iterations, each container
	has size at most $(1-\delta)^t\binom{n}{2}$ yielding a maximum of
	$(1-\delta)^t n^3$ triangles.
	So we must have that $(1-\delta)^t \geq \eps$, and therefore $t \leq
	\log(1/\eps)/\delta$. So how many containers do we end up with? Something like
	\[
		\binom{n^2}{n^{3/2}}^t = n^{O(n^{3/2})} = \exp(O(n^{3/2}\log n)).
	\]
	We can even bound the sizes of the fingerprints by $t \tau
	\binom{n}{2} = Cn^{3/2}$.
	% Let $\mathcal{C}_0 = \{\mathcal{H}_{K_3}(K_n)\}$. What
\end{proof}

\subsection{Some cool applications!}

\subsubsection{Counting triangle-free graphs}

We will try to prove the following theorem:
\begin{theorem}
	\label{containers:far_bipartite}
	Almost all triangle free graphs with $n$ vtxs and $m \gg n^{3/2}$
	edges are $o(m)$ close to being bipartite.
\end{theorem}

To prove it we will need an estability result (basically
Erdös-Simonitivs) that states
\begin{theorem}
	If $G$ is $t$-far from being bipartite, then
	\[
		k_3(G) \geq \frac{n}{6}(e(G) + t - \frac{n^2}{4})
	\]
\end{theorem}

\begin{proof}
	(of \ref{containers:far_bipartite}) Let $X(n,m,\alpha)$ be the
	number  of  triangle free graphs with $m$ edges which are $\alpha
	m$-far from being bipartite.
	We will upperbound $X$ counting by counting inside the containers.
	So suppose we applied [\ref{trm:container_for_triangles}] with
	$\eps>0$, which yielded $\mathcal{G}$.
	We have
	\[
		X(n,m,\alpha) \leq \sum_{G \in \mathcal{G}} \#\{ K_3 \not \subset H
		\subset G,\; e(H) = m,\; H\,\alpha m\text{-far from being bip.}\}.
	\]
	Now fix some parameter $\delta > 7\eps$ we choose later, we know
	from the stability result that either $G$ is $\delta n^2$ close to
	being bipartite or
	\[
		e(G) \leq \frac{n^2}{4} + 6\eps n^2 - \delta n^2 \leq \frac{n^2}{4}
		- \eps n^2.
	\]
	So let's separate our counting depending on the container and bound
	the previous sum
	\[
		X(n,m,\alpha) \leq \sum_{\substack{G \in \mathcal{G}\\e(G) \leq
		(1/4 - \eps)n^2}} \binom{e(G)}{m}\quad
		+ \quad\sum_{\substack{G \in \mathcal{G}\\G\, \delta-\text{close to
		bip.}}} \binom{\delta n^2}{\alpha m} \binom{n^2/4}{(1-\alpha)m}
	\]
	where the first inequality follows from taking all $m$ subgraphs of
	$G$ when $G$ has few edges and the second follows from choosing the
	$\alpha m$ edges
	that make $H$ far from being bipartite. A bit of manipulation shows
	\begin{align*}
		X(n,m,\alpha) &\leq \sum_{G \in \mathcal{G}} (1-4\eps)^m
		\binom{n^2/4}{m} + \binom{\delta n^2}{\alpha
		m}\binom{n^2/4}{m}\bigg(\frac{n^2}{4}\bigg)^{-\alpha m}\\
		&\leq n^{Cn^3/2}e^{-cm} \binom{n^2/4}{m} \to o(1)\binom{n^2/4}{m}
	\end{align*}
	when $m \gg n^{3/2}\log n$.

	What if we want to remove the $\log n$? Then we do the same thing
	but we use our fingerprints, we sum over them and remember that each
	$S$ must be cointained in $H$.
	Taking that into consideration the $\log n$ vanishes.

	\begin{align*}
		X(n,m,\alpha) \leq \sum_{|S| = 0}^{Cn^{3/2}}
		\binom{n^2/2}{|S|}\bigg(\binom{(1/4-\eps)n^2}{m - |S|} +
		\binom{n^2/4}{(1-\alpha/2)m - |S|}\bigg)
	\end{align*}
	and after some more careful computation we win.
\end{proof}

\subsubsection{When does Gnp arrows a triangle?}

The following is a very surprising theorem (and beautiful one as well).
\begin{theorem}
	\label{containers:arrows_k3}
	Let fix $r \in \N$ and let $p \gg n^{-1/2}$, then $G(n,p) \to_r K_3$ w.h.p.
\end{theorem}
As always, we need some kind of saturation result. We will actually
show that $\Pr(G(n,p) \not \to K_3) \to 0$. So what does this event really mean?
If $G(n,p) \not \to_r K_3$, then there exists a partition (or
coloring) of the edges of $G(n,p) = H_1 \cup H_2 \dots \cup H_r$,
such that each $H_i$ is $K_3$-free.
Now an easy double counting yields the following lemma:
\begin{lemma}
	Let $G$ be a graph and $\chi$ an r-coloring of its edges. Then if
	\[
		e(G) \geq (1-\delta)\binom{n}{2}
	\]
	then $\chi$ contains $\delta n^3$ monochromatic triangles.
\end{lemma}
And we are ready to apply our containers.
\begin{proof}
	(of \ref{containers:arrows_k3}) We upperbound again on our
	containers, let $C_i$ be the containers for $H_i$. We know $E(K_n) =
	F \cup C_1 \dots \cup C_r$
	where $k_3(C_i) \leq \eps n^3$ for each $i$ and $F$ is the
	complement of the union.
	By the previous lemma, it must be that $e(F) \geq \eps r
	\binom{n}{2}$, and this means that, if $G(n,p) \subset C_1 \cup
	\dots \cup C_r$,
	it is avoiding a positive density set of edges. We find that
	(maximizing by the last term)
	\begin{align*}	
		\Pr(G(n,p) \not \to K_3) &\leq \sum_{S_1 \cup \dots \cup S_r}
		\Pr((S_1 \cup \dots S_r) \subset G(n,p) \subset (C_1 \cup \dots \cup C_r))\\
		&\leq \sum_{t = 0}^{Crn^{3/2}} \binom{n^2/2}{t}p^t(1-p)^{\delta
		n^2} \leq e^{-p\delta n^2}\bigg(\frac{epn^2}{Crn^{3/2}}\bigg)^{Crn^{3/2}}\\
		&= \bigg(\frac{e^{-\delta p\sqrt{n}}p\sqrt{n}}{D}\bigg)^{Crn^{3/2}}
	\end{align*}
	and it is clear that if $p\sqrt{n} \to \infty$, then the probability
	is tends to 0.
\end{proof}

\subsection{A modern approach to Containers}

We've seen some simple applications of the Hypergraph Container Lemma
[\ref{lemma:hypergraph_containers}] but no proof of it. In this section,
our objective will be to develop a new and modern proof (beyond the
Kleitman-Winston algorithm).

\subsubsection{Graph Containers}
The following theorem is similar to what you get when you apply
[\ref{lemma:hypergraph_containers}] for graphs.
\begin{theorem} [Graph Containers]
	\label{trm:graph_containers}
	Let $0 < p < \alpha < 1$, given a graph $G$ with vtxs set $V$ and
	$|V| = n$, there exists a family of containers
	$\mathcal{C} \subset 2^V$, satisfying the following properties:
	\begin{enumerate}[label=(\roman*)]
		\item $|\mathcal{C}| \leq \binom{n}{pn/\alpha}$
		\item for all $I \in \mathcal{I}(G)$, there exists $C \in
			\mathcal{C}$ s.t $I \subset C$
		\item for every $C \in \mathcal{C}$, $e(G[C]) \leq \frac{4\alpha|C|}{p}$
	\end{enumerate}
\end{theorem}

We know the usual (naive) way to prove this.
\begin{proof}[Sketch of proof]
	For any independent set $I$, we start off with an empty fingerprint
	set $S = \varnothing$ and the whole vtx set $V$ as $C$.
	We iteratively query each vertex in the highest degree inside $C$
	order if it belongs to $I$, if so we add it to $S$ and remove its
	neighbors from $C$,
	if not we remove it from $C$. A simple counting then yields
	properties similar to the ones listed if we stop $S$ at size at most
	$pn/\alpha$.
\end{proof}

Marcelos's approach will be very similar to this (although
considerably more complicated), we will still query the vtxs of the graph
in a smart order, but we will use the best information to do so, to
build the best containers! We will add vertices that most reduce the space
of independent sets they belong to. In some sense, the best possible
fingerprints you could hope for if you chose randomly.

Instead of the maximal degree in the remaining world, we will
consider the following magic statistic of a vertex $v$
\[
	\Pr(v \in C_p \,| \,C_p \in \mathcal{I}(G[C])),
\]
where $C_p$ is a $p$-random subset of $C$, which states roughly in
how many independent sets the vertex $v$ belongs to in $G[C]$. If it
belongs to few, it is
a good fingerprint candidate. Note that by Harris, as one event is
increasing, and the other decreasing,
\[
	\Pr(v \in C_p \,| \,C_p \in \mathcal{I}(G[C])) = \frac{\Pr(v \in C_p
	\land C_p \in \mathcal{I}(G[C]))}{\Pr(C_p \in \mathcal{I}(G[C]))}
	\leq \Pr(v \in C_p) = p.
\]

This motivates the following algorithm, given $I \in \mathcal{I}(G)$:
\begin{enumerate}
	\item Initialize $S = \varnothing, C = V$.
	\item While $\exists v \in C$ s.t $\Pr(v \in C_p\, |\, C_p \in
		\mathcal{I}(G[C])) < (1-\alpha)p$ (pick $v$ minimal)
		\begin{enumerate}
			\item If $v \in I$: $S \leftarrow S \cup \{v\}$, $C \leftarrow C
				\setminus N(v)$.
			\item If $v \not \in I$: $C \leftarrow C \setminus \{v\}$.
		\end{enumerate}
\end{enumerate}

We need to check that in the end, this algorithm yields what we want,
a deterministic fingerprint $S$ with $|S| < pn/\alpha$
and a container associated with it with few edges. It is clear that,
by design, $I \subset C$.

This proof is a bit magical, we smartly defined this algorithm, now
let's keep track of the probability of the event that it is tracking throughout.
The probability bounds we get from this event will yield properties
(i) and (iii) from [\ref{trm:graph_containers}].

Suppose then, that given $I$, the algorithm queries the vtxs $v_1,
\dots, v_M$. We want to know
that, given that we chose $V_p$ randomly from the independent sets,
what is the probability we got our fingerprint.
To that end, define the events
\begin{align}
	\mathcal{E}_0 &= \{V_p \in \mathcal{I}(G)\}\\
	\mathcal{E}_i &= \mathcal{E}_{i-1} \cap
	\begin{cases}
		\{v_i \in V_p\} & \text{ if } v_i \in S\\
		\{v_i \not \in V_p\} & \text{ if } v_i \not \in S\\
	\end{cases}
\end{align}
Smartly calculating $\Pr(\mathcal{E}_i)$ will yield all properties :).
But first we must see this event in
a way more recognizable by our algorithm. Let $S^i$ and $C^i$
represent the fingerprint and the container after querying the $i$th vtx,
we have

\begin{claim}
\[ \mathcal{E}_i = \{V_p \in \mathcal{I}(G[C^{i}])\} \cap \{V_p
\subset C^{i}\} \cap \{S^i \subset V_p\} \]
\end{claim}

\begin{proof}
To prove the inclusion (lhs to rhs), it suffices to notice that,
during the algorithm, we remove precisely only vtxs that got denied
and neighbors of vtxs that got accepted into the fingerprint (so if
	$V_p \in \mathcal{I}(G)$, it must be contained in $C^i$, and
therefore independent in $G[C^i]$).
Also it is clear that $\mathcal{E}_i \subset \{S^i \subset V_p\}$.
The other inclusion is basically the same.
\end{proof}
In particular, we find that
\[
\Pr(\mathcal{E}_i) = \Pr(v_i \in V_p \land \mathcal{E}_{i-1}) =
\Pr(v_i \in V_p | \mathcal{E}_{i-1}) \Pr(\mathcal{E}_{i-1})
\]
And that motivates one more claim (to be used in the algorithm)
\begin{claim}
\[
	\Pr(v_i \in V_p | \mathcal{E}_{i-1}) = \Pr(v_i \in C^{i-1}_p |
	C^{i-1}_p \in \mathcal{I}(G)) < (1-\alpha)p \Pr(\mathcal{E}_{i-1})
\]
where here $v_i$ is such that $\Pr(v \in C^{i-1}_p\, |\, C^{i-1}_p
\in \mathcal{I}(G)) < (1-\alpha)p$
\end{claim}

Note that if this were true, by the way we choose the vtxs in the
main loop, we would find that $\Pr(\mathcal{E}_M) < (1-\alpha)^{|S|}p^{|S|}$

\begin{proof}[Proof of Claim]
We want to show
\[
	\Pr(v_i \in V_p | V_p \in \mathcal{I}(G[C^{i-1}]) \land V_p
	\subset C^{i-1} \land S^{i-1} \subset V_p) = \Pr(v_i \in C^{i-1}_p
	| C^{i-1}_p \in \mathcal{I}(G[C^{i-1}]))
\]

The last equality is clear from the fact that we are sampling only
from $C$, the first one follows from
the fact that conditioning on $\{V_p \subset C^{i-1}\}$ and
$\{V_p \in \mathcal{I}(G)\}$, the event  $\{S^{i-1} \subset V_p\}$ is
independent from $\{v_i \in
V_p\}$ (as $S^{i-1} \cup \{v_i\}$ is an indep. set).
\end{proof}

This already gives the first property!
\begin{claim}
(Property (i) of [\ref{trm:graph_containers}]) At the end of the
algorithm, the fingerprint produced has size at most $pn/\alpha$.
\end{claim}
\begin{proof}
Suppose as before, that given $I$, we queried the vtxs $v_1, \dots,
v_M$ and the algorithm stopped producing
a fingerprint $S$.
We note that
\[
	p^{|S|}(1-p)^{n - |S|} = \Pr(V_p = S) \leq \Pr(\mathcal{E}_M) \leq
	(1-\alpha)^{|S|}p^{|S|}.
\]
In particular, approximating $(1-p)$ by $e^{-p}$ and $(1-\alpha)$ by
$e^{-\alpha}$,
we find that
\[e^{p(n - |S|)} \leq e^{-\alpha|S|}\]
so that $|S| \leq pn/\alpha$.
\end{proof}
